\documentclass[10pt,a4paper]{amsart}
\usepackage[margin=19mm]{geometry}
\usepackage{amsmath,amssymb,mathtools}
\usepackage[scr=rsfs,cal=euler]{mathalfa}
\usepackage{xfrac}
\usepackage[unicode,hidelinks,bookmarks=false]{hyperref}
\newcommand{\ie}{i.e.\ }
\newcommand{\cf}{cf.\ }
\newcommand{\resp}{respectively\ }
\newcommand{\Subsection}{\subsection}
\providecommand{\nobreakdash}{\nobreak\hbox{-}}
\setlength{\emergencystretch}{2em}
\multlinegap0pt
\usepackage{amssymb}
\newtheorem{lemma}{Lemma}
\DeclareMathOperator{\di}{d}
\title{Shadowing and Hyperbolicity for Endomorphisms of Locally Compact Groups}
\author{Dekui Peng}
\date{Source: arXiv:2606.27647v3 (CC BY 4.0)}
\begin{document}
\maketitle

\noindent\emph{Source and licence.} Shadowing and Hyperbolicity for Endomorphisms of Locally Compact Groups. Version arXiv:2606.27647v3 (CC BY 4.0), distributed by arXiv under the Creative Commons Attribution 4.0 International licence, http://creativecommons.org/licenses/by/4.0/, as recorded in the arXiv licence metadata for this exact version and shown on the abstract page https://arxiv.org/abs/2606.27647v3. Full licence notice: \texttt{licenses/arxiv-2606.27647-CC-BY-4.0.txt}.\par\smallskip

\noindent\textit{Editorial scope.} Counted unit: the article's lemma lem:local-error-equation (source0.tex lines 622--632). The article's complete original proof of it is delivered with it (source0.tex lines 634--844, one \texttt{proof} environment of 211 lines). No mathematics is written, repaired or re-derived: the statement and the proof are the article's own lines, in the article's own notation, and no retained line is substituted or re-flowed.  No further source-local context is needed: every cross-reference of every retained line resolves inside the two delivered spans.  Nothing was omitted from any retained span, so the package carries no omission record.  No retained line contains a citation, so the wrapper carries no bibliography and no bibliography entry.  No retained line contains a figure environment, an image-inclusion command or drawing code, so no image is delivered and the package declares no asset dependency.  Editorial formatting only, in addition: an AMS \texttt{amsart} preamble that restates the article's own declarations and macros used by the retained lines, namely the environments \texttt{lemma} on separate counters, together with the standard packages those lines need.  The retained environments print in this document as Lemma 1 in order of appearance, whereas the article prints the same items with the numbering of its own full document.  The complete native source of the article as distributed in the arXiv e-print for this exact version is bundled unchanged as \texttt{sources/203785/source0.tex}.
\begin{lemma}\label{lem:local-error-equation}
Let $G$ be a Lie group with Lie algebra $\mathfrak g$, let $\alpha:G\to G$ be a continuous endomorphism,
and put $A=\di\alpha:\mathfrak g\to\mathfrak g$. If $A$ is hyperbolic, then
for every neighbourhood $U$ of $e$ in $G$, there exists a neighbourhood
$V$ of $e$ such that, for every sequence $(e_n)_{n\geq 0}$ in $V$, the
equation
\begin{equation}\label{eq:error-equation}
d_{n+1}=\alpha(d_n)e_n,\qquad n\geq 0,
\end{equation}
has a solution $(d_n)_{n\geq 0}$ contained in $U$.
\end{lemma}

\begin{proof}
Choose an exponential chart $\theta:O\to E$ from a neighbourhood $O$ of
$e$ in $G$ onto a neighbourhood of $0$ in $E=\mathfrak g$, with
$\theta(e)=0$ and $(D\theta)(e)=\operatorname{id}_{\mathfrak g}$. We shall
shrink $O$ several times below.

In these coordinates, put
\[
\mathcal D=\{(\xi,a)\in\theta(O)\times\theta(O):
\alpha(\theta^{-1}(\xi))\theta^{-1}(a)\in O\}.
\]
Then $\mathcal D$ is an open neighbourhood of $(0,0)$. Define the smooth
map $F:\mathcal D\to E$ by
\[
F(\xi,a)=\theta\bigl(\alpha(\theta^{-1}(\xi))\theta^{-1}(a)\bigr),
\]
so that $F(0,0)=0$ and $D_\xi F(0,0)=A$. Write
\[
F(\xi,a)=A\xi+\psi(\xi,a).
\]
Thus $\psi(0,0)=0$ and $D_\xi\psi(0,0)=0$.

Since $A$ is hyperbolic, there is an $A$-invariant decomposition
\begin{equation}\label{eq:splitting}
E=E^{\mathrm{s}}\oplus E^{\mathrm{u}},
\end{equation}
where the spectrum of $A_{\mathrm{s}}=A|_{E^{\mathrm{s}}}$ lies inside the
open unit disk in $\mathbb{C}$, and $A_{\mathrm{u}}=A|_{E^{\mathrm{u}}}$ is invertible with
spectrum outside the closed unit disk. Choose norms on
$E^{\mathrm{s}}$ and $E^{\mathrm{u}}$, and use the product norm
\[
\|\xi\|=\max\{\|\xi^{\mathrm{s}}\|,\|\xi^{\mathrm{u}}\|\}
\]
on $E$, so that, for some $0<\lambda<1$,
\begin{equation}\label{eq:lambda-estimate}
\|A_{\mathrm{s}}\|\leq\lambda,\qquad
\|A_{\mathrm{u}}^{-1}\|\leq\lambda.
\end{equation}

We next choose the constants used in the contraction argument. Since
$D_\xi\psi(0,0)=0$ and $D_\xi\psi$ is continuous as a function of the pair
$(\xi,a)$, we may make the Lipschitz constant of $\psi$ in the first
variable arbitrarily small, uniformly for all sufficiently small values of
the second variable. More precisely, choose $L>0$ with
\begin{equation}\label{eq:L-small}
\frac{L}{1-\lambda}<\frac14.
\end{equation}
Then choose $\rho>0$ so small that
\[
\overline{B_\rho(0)}\times\overline{B_\rho(0)}\subseteq\mathcal D,
\qquad
\theta^{-1}(\overline{B_\rho(0)})\subseteq U,
\]
and
\[
\|D_\xi\psi(\xi,a)\|\leq L
\]
whenever $\|\xi\|\leq\rho$ and $\|a\|\leq\rho$. Since the closed ball is
convex, the mean value estimate gives, for
$\|\xi\|,\|\zeta\|\leq\rho$ and $\|a\|\leq\rho$,
\begin{equation}\label{eq:psi-lipschitz}
\|\psi(\xi,a)-\psi(\zeta,a)\|\leq L\|\xi-\zeta\|.
\end{equation}

Since $\psi(0,0)=0$, choose $0<\eta<\rho$ such that
\begin{equation}\label{eq:psi-zero-small}
\|\psi(0,a)\|\leq \frac{(1-\lambda)\rho}{4}
\end{equation}
whenever $\|a\|\leq\eta$. Combining \eqref{eq:L-small},
\eqref{eq:psi-lipschitz},  and
\eqref{eq:psi-zero-small}, we obtain, for all $\|\xi\|\leq\rho$ and
$\|a\|\leq\eta$,
\begin{equation}\label{eq:psi-bound}
\|\psi(\xi,a)\|
\leq
\|\psi(\xi,a)-\psi(0,a)\|+\|\psi(0,a)\|
\leq
L\rho+\frac{(1-\lambda)\rho}{4}
\leq
\frac{(1-\lambda)\rho}{2}.
\end{equation}

Let $V=\theta^{-1}(B_\eta(0))$. Now let $(e_n)_{n\geq 0}$ be a sequence
in $V$, and put $a_n=\theta(e_n)$. We shall solve
\begin{equation}\label{eq:coordinate-recurrence}
\xi_{n+1}=A\xi_n+\psi(\xi_n,a_n)
\end{equation}
with $\|\xi_n\|\leq\rho$ for all $n\geq 0$. Then
$d_n=\theta^{-1}(\xi_n)$ will solve \eqref{eq:error-equation} and will
belong to $U$.

Let $\mathcal X$ be the complete metric space of all bounded sequences
$\xi=(\xi_n)_{n\geq 0}$ in $E$ satisfying $\|\xi\|_\infty\leq\rho$. For
$\xi\in\mathcal X$, write
\[
\psi(\xi_n,a_n)
=
\psi^{\mathrm{s}}(\xi_n,a_n)+\psi^{\mathrm{u}}(\xi_n,a_n)
\]
according to the decomposition \eqref{eq:splitting}. Define
$\Gamma:\mathcal X\to \ell^\infty(E)$ by
\begin{align}
(\Gamma\xi)_n^{\mathrm{s}}
&=
\sum_{k=0}^{n-1}A_{\mathrm{s}}^{\,n-1-k}
\psi^{\mathrm{s}}(\xi_k,a_k),
\label{eq:Gamma-stable}\\
(\Gamma\xi)_n^{\mathrm{u}}
&=
-\sum_{k=n}^{\infty}A_{\mathrm{u}}^{\,n-1-k}
\psi^{\mathrm{u}}(\xi_k,a_k).
\label{eq:Gamma-unstable}
\end{align}
The second series converges because
$\|A_{\mathrm{u}}^{-1}\|\leq\lambda<1$.

We now check that $\Gamma$ maps $\mathcal X$ into itself. Put
\[
M_\xi=\sup_{n\geq 0}\|\psi(\xi_n,a_n)\|.
\]
By \eqref{eq:psi-bound}, $M_\xi\leq (1-\lambda)\rho/2$. For the stable
component, using \eqref{eq:lambda-estimate}, we have
\[
\|(\Gamma\xi)_n^{\mathrm{s}}\|
\leq
\sum_{k=0}^{n-1}\lambda^{n-1-k}M_\xi
\leq
\frac{M_\xi}{1-\lambda}.
\]
For the unstable component,
\[
\|(\Gamma\xi)_n^{\mathrm{u}}\|
\leq
\sum_{k=n}^{\infty}\lambda^{k+1-n}M_\xi
\leq
\frac{M_\xi}{1-\lambda}.
\]
Since the norm on $E$ is the maximum of the stable and unstable norms, it
follows that
\[
\|\Gamma\xi\|_\infty
\leq
\frac{M_\xi}{1-\lambda}
\leq
\frac{\rho}{2}
<\rho.
\]
Thus $\Gamma(\mathcal X)\subseteq\mathcal X$.

Next we show that $\Gamma$ is a contraction. Let $\xi,\zeta\in\mathcal X$
and put
\[
M_{\xi,\zeta}
=
\sup_{n\geq 0}
\|\psi(\xi_n,a_n)-\psi(\zeta_n,a_n)\|.
\]
By \eqref{eq:psi-lipschitz},
$M_{\xi,\zeta}\leq L\|\xi-\zeta\|_\infty$. Again, the estimates in \eqref{eq:lambda-estimate} give
\[
\|(\Gamma\xi-\Gamma\zeta)_n^{\mathrm{s}}\|
\leq
\frac{M_{\xi,\zeta}}{1-\lambda},
\qquad
\|(\Gamma\xi-\Gamma\zeta)_n^{\mathrm{u}}\|
\leq
\frac{M_{\xi,\zeta}}{1-\lambda}.
\]
Therefore,
\[
\|\Gamma\xi-\Gamma\zeta\|_\infty
\leq
\frac{L}{1-\lambda}\|\xi-\zeta\|_\infty
<
\frac14\|\xi-\zeta\|_\infty.
\]
Hence $\Gamma$ is a contraction on $\mathcal X$. By the Banach fixed point
theorem, $\Gamma$ has a fixed point $\xi=(\xi_n)_{n\geq 0}$ in $\mathcal X$.

It remains to check that this fixed point solves
\eqref{eq:coordinate-recurrence}. For the stable component, the definition
\eqref{eq:Gamma-stable} gives
\[
\xi_{n+1}^{\mathrm{s}}
=
A_{\mathrm{s}}\xi_n^{\mathrm{s}}
+
\psi^{\mathrm{s}}(\xi_n,a_n).
\]
For the unstable component, shifting the series in
\eqref{eq:Gamma-unstable} gives
\[
\xi_{n+1}^{\mathrm{u}}
=
A_{\mathrm{u}}\xi_n^{\mathrm{u}}
+
\psi^{\mathrm{u}}(\xi_n,a_n).
\]
Thus
\[
\xi_{n+1}=A\xi_n+\psi(\xi_n,a_n)=F(\xi_n,a_n).
\]
Consequently, with $d_n=\theta^{-1}(\xi_n)$, we have
\[
d_{n+1}=\alpha(d_n)e_n
\]
for all $n\geq 0$. Since $\|\xi_n\|\leq\rho$ and
$\theta^{-1}(\overline{B_\rho(0)})\subseteq U$, the sequence
$(d_n)_{n\geq 0}$ is
contained in $U$. This completes the proof.
\end{proof}

\end{document}
